\subsection{复对数}

设 B 是由点 \(z = x + iy\) （满足 \(-\pi \leqslant y < \pi\) ）所成的“带形”；函数 \(e^{z}\) 在 B 上取它的每个值一次且仅一次；换句话说，\(z \mapsto e^{z}\) 是 B 到 \(C^{*}\) 上的连续双射；在此映射下，（半开）线段 \(x = x_{0}\) ，\(-\pi \leqslant y < \pi\) 的像是圆 \(|z| = e^{x_{0}}\) ；直线 \(y = y_{0}\) 的像是由 \(\operatorname{Am}(z) = y_{0} \pmod{2\pi}\) 定义的（开）半直线。在映射 \(z \mapsto e^{z}\) 下，B 的内部 B （即 \(z \in C\) 中满足 \(|\mathcal{I}(z)| < \pi\) 者所成的集）的像是（闭的）负实半轴在 C 中的余集 F；若约定用 \(\operatorname{Am}(z)\) 表示属于 \([-\pi, \pi]\) 的 z 的辐角量度，则集 F 可由关系 \(-\pi < \operatorname{Am}(z) < \pi\) 定义。由于 \(z \mapsto e^{z}\) 是 C 到 \(C^{*}\) 上的严格同态，任一含于 B 的开集（从而含于 C）在此映射下的像是 \(C^{*}\) 中（从而是 F 中）的开集；换句话说，\(z \mapsto e^{z}\) 在 B 上的限制是 B 到 F 上的同胚。我们把后者的逆同胚，即 F 到 B 上的同胚记作 \(z \mapsto \log z\) ；对复数 \(z \in F\) ，\(\log z\) 称为 z 的对数的主值。若 \(z = x + iy\) 且 \(\log z = u + iv\) ，则 \(x + iy = e^{u + iv}\) ，由此 \(e^{u} = |z|\) ，并且由于 \(-\pi < v < \pi\) ，有 \(v = \operatorname{Am}(z)\) 。此外，有 \(\tan(v + \pi/2) = -x/y\) （若 \(y \neq 0\) ）；于是可以写出

\[
\left\{ \begin{array}{l l} u = \log | z | = \frac {1}{2} \log (x ^ {2} + y ^ {2}) \\ v = \frac {\pi}{2} - \operatorname{Arc} \tan \frac {x}{y} & \text { if } y > 0 \\ v = 0 & \text { if } y = 0 \\ v = - \frac {\pi}{2} - \operatorname{Arc} \tan \frac {x}{y} & \text { if } y <   0. \end{array} \right.\tag{30}
\]

显然 \(\log z\) 是函数 \(\log x\) （定义在正开实半轴 \(R_{+}^{*}\) 上）到 F 上的延拓。若 z，\(z'\) 是 F 中使 \(zz'\) 非负实数的两点，则有 \(\log(zz') = \log z + \log z' + 2\varepsilon\pi i\) ，其中 \(\varepsilon = +1\) 、-1 或 0，视 \(\mathrm{Am}(z)\) 与 \(\mathrm{Am}(z')\) 的值而定。

我们注意，在负实半轴的各点处函数 \(\log z\) 没有极限；确切地说，若 x 趋于 \(x_{0} < 0\) 且 y 保持 \textgreater{} 0（相应地，\textless{} 0）地趋于 0，则 \(\log z\) 趋于 \(\log |x_{0}| + \pi i\) （相应地，\(\log |x_{0}| - \pi i\) ）；当 z 趋于 0 时，\(|\log z|\) 无限增大。

以后我们将看到，解析函数理论如何使我们能把函数 \(\log z\) 延拓，并在完全一般的意义下定义复对数。

由于 \(\log z\) 是 \(e^{z}\) 的逆同胚，反函数求导公式（I, p.~9, prop. 6）表明 \(\log z\) 在每一点 \(z \in F\) 处可导，并且

\[
\mathrm{D} (\log z) = \frac {1}{e ^ {\log z}} = \frac {1}{z}\tag{31}
\]

这个公式推广了 III, p.~93 的公式 (10)。

